\section{Introduction}\label{sec:intro}

\subsection{Background and goal}

The module covers three themes: GI science, GI systems and 3D GIS. Instead of one large project, I chose to make one small web application for each theme. Each project starts from a single question (\cref{tab:overview}), and I wanted every app to be finished and usable by others, not only by me. The three apps are therefore published together on one website, \siteurl, where a homepage links to each of them (\cref{fig:mobile} in \cref{sec:website}). The source code is in the public repository \url{https://github.com/progrimm/TBA4565-Geomatics-Specialization-Course}.

\begin{table}[htbp]
\centering
\caption{The three projects.}
\label{tab:overview}
\small
\begin{tabularx}{\textwidth}{L{2.6cm} L{1.7cm} Y}
\toprule
Project & Theme & Question \\
\midrule
Projection Lab & GI science & How much does the Mercator projection inflate the area of each country, and how does this depend on latitude? \\
\addlinespace
Trondheim 3D & GI systems & Where is the nearest public shelter in Trondheim, and how far is the walk? \\
\addlinespace
3DGS Viewer & 3D GIS & What does a 3D Gaussian splat file contain, and can it be opened and inspected in a browser without a server? \\
\bottomrule
\end{tabularx}
\end{table}

\subsection{Way of working}\label{sec:work}

I used the AI coding assistant Claude Code \citep{claudecode} in all three projects, in two different roles. In the first part of Projection Lab, it worked as a tutor: I wrote and ran the code myself, and the assistant explained concepts and pointed me to the official documentation. In the rest of the work, the assistant wrote the code, while I set the goals, chose between alternatives, reviewed the changes and checked the result in the browser.

I chose this way of working on purpose. AI assistants are becoming a normal part of software development, and I expect that working in geomatics will involve directing and reviewing such tools as much as writing code by hand. I wanted to experience GIS development from this perspective, while still learning the underlying technologies well enough to understand and check what the assistant made.

\section{Data and tools}\label{sec:data}

All data used in the projects is open, except my own splat capture in \cref{sec:p3}. \Cref{tab:data} lists the datasets, and all screenshots in this report contain data from these sources. \Cref{tab:tools} lists the main tools and libraries, and what each one is used for.

\begin{table}[!ht]
\centering
\caption{Datasets.}
\label{tab:data}
\small
\begin{tabularx}{\textwidth}{L{3.6cm} L{2.6cm} L{1.9cm} Y}
\toprule
Dataset & Source & Licence & Description \\
\midrule
Admin 0 Countries, 1:110m & Natural Earth \citep{naturalearth} & Public domain & Boundaries of 177 countries (GeoJSON, WGS 84). Used in Projection Lab. \\
\addlinespace
Public shelters (\emph{Tilfluktsrom -- offentlige}) & DSB, WFS service \citep{dsbwfs} & NLOD 2.0 & 12 shelters with number, capacity and address, in EUREF89 / UTM zone 32N. Used in Trondheim 3D. \\
\addlinespace
Walking network & OpenStreetMap, Overpass API \citep{overpass} & ODbL & Footways, paths and streets around the shelters: 164\,294 nodes and 32\,944 ways. Used in Trondheim 3D. \\
\addlinespace
Cesium World Terrain and Cesium OSM Buildings & Cesium ion & Cesium ion terms; buildings ODbL & Global terrain and OpenStreetMap buildings, streamed as needed. Used in Trondheim 3D. \\
\addlinespace
Splat files & Own capture; ``SOAR Telescope'' by L.~Cerda \citep{soar} & Not published; CC BY 4.0 & Test files for the 3DGS Viewer. The app itself ships no data. \\
\bottomrule
\end{tabularx}
\end{table}

\begin{table}[!ht]
\centering
\caption{Main tools and libraries.}
\label{tab:tools}
\small
\begin{tabularx}{\textwidth}{L{3.6cm} L{2.6cm} Y}
\toprule
Tool & Project & Function \\
\midrule
D3 (d3-geo) \citep{d3} & Projection Lab & Projects and draws the world map in the browser \\
PostgreSQL with PostGIS & Projection Lab & Stores the countries; geodesic area with \code{ST\_Area} \citep{postgis-starea} \\
GDAL ogr2ogr \citep{gdal-ogr2ogr} & Projection Lab & Imports the GeoJSON file into PostGIS \\
pyproj \citep{pyproj} and Shapely \citep{shapely} & Projection Lab & Spherical, Mercator and Eckert IV areas \\
FastAPI \citep{fastapi} & Projection Lab & Web API between the database and the map \\
Docker Compose & Projection Lab & Runs the database, the API and the web server together \\
CesiumJS \citep{cesiumjs} and Cesium ion & Trondheim 3D & 3D globe, terrain and streaming of 3D Tiles \\
proj4js & Trondheim 3D & Converts the shelters from EPSG:25832 to WGS 84 \\
Spark \citep{spark} and three.js \citep{threejs} & 3DGS Viewer & Renders Gaussian splats with WebGL \\
React, Vite & All & User interface and build \\
Vitest, pytest, ESLint & All & Unit tests and code checks \\
GitHub Actions and Pages \citep{ghpages} & All & Automatic checks and hosting of the website \\
\bottomrule
\end{tabularx}
\end{table}
